
\section{Komplexe Inhaltselemente}
\subsection{Punktliste}

\begin{itemize}
    \item Wird
    \item erstellt
    \item mit
    \item itemize
\end{itemize}

\subsection{Geordnete Liste}

\begin{enumerate}
    \item Wird
    \item erstellt
    \item mit
    \item enumerate
\end{enumerate}

\subsection{Definitionsliste}

\begin{description}
    \item [Was beschrieben] Wird
    \item [werden soll] erstellt
    \item [kommt in] mit
    \item [eckige Klammern] description
\end{description}

\subsection{Geschachtelte Liste}

\begin{enumerate}
    \item hier
    \item wird
    \begin{itemize}
        \item eine Liste
        \item innerhalb einer Liste
    \end{itemize}
    \item angelegt.
\end{enumerate}

\subsection{Tabellen}

LaTeX platziert Tabellen eigenständig irgendwo.
Man kann versuchen LaTeX Hinweise zu geben, wo die
Tabellen hinsollen mit "`htb"': here, top, bottom (der Seite).

\begin{table}[h]
    \caption{Dies ist die Caption einer Tabelle}
    % eigentlichen Tabelleninhalte ab hier:
    % die cols werden beschrieben in den zweiten geschweiften klammern
    % | steht für den strich der tabelle, lcr für die Textbündigkeit links, center, rechts
    % die horizontalen Lininen sind \hline
    % auch in Tabellen Zeilenende mit \\ nötig
    \begin{tabular}{|l|c|r|}
        \hline
        \textbf{Überschrift-fett} & \textsc{Überschrift-Kapitälchen} & \texttt{Überschrift-Schreibmaschine} \\
        \hline
        Dies ist & Dies ist & Dies ist \\
        \hline
        linksbündig & zentriert & rechtsbündig \\
        \hline
    \end{tabular}
\end{table}

\noindent Tipp: Tabelle mit Excel erstellen und LaTeX-code ausgeben lassen.




\subsection{Abbildungen}











\subsection{Mathematische Formeln}
\subsection{URLs}